Respuesta oral breve, con un bosquejo a mano. Considere, para \(t\ge0\),
\[
\mathrm A:\ y''+2y'=1,
\qquad
\mathrm B:\ y''+2y=1,
\qquad y(0)=y'(0)=0\ \text{en ambos casos}.
\]
Clasifique las ecuaciones y elija un primer paso útil para cada una, justificando qué datos emplea y cómo recupera los restantes. Sin desarrollar una resolución completa de ambas, dibuje la respuesta A: indique la tangente inicial, el crecimiento, la concavidad y el comportamiento de la pendiente cuando el tiempo aumenta. Explique por qué la semejanza de los enunciados no permite trasladar ese bosquejo a B.
